% Generated draft: requires mathematical and source review.
\documentclass[11pt]{article}
\usepackage[margin=25mm]{geometry}
\usepackage{fontspec}
\setmainfont{Latin Modern Roman}
\usepackage{amsmath,amssymb,mathtools,unicode-math}
\setmathfont{Latin Modern Math}
\usepackage{enumitem,needspace,xurl,hyperref}
\setlength{\parindent}{0pt}
\setlength{\parskip}{5pt}
\setlength{\emergencystretch}{3em}
\newcommand{\N}{\mathbb{N}}
\newcommand{\Z}{\mathbb{Z}}
\newcommand{\Q}{\mathbb{Q}}
\newcommand{\R}{\mathbb{R}}
\newcommand{\C}{\mathbb{C}}
\newcommand{\F}{\mathbb{F}}
\newcommand{\A}{\mathbb{A}}
\newcommand{\PP}{\mathbb{P}}
\newcommand{\E}{\mathbb{E}}
\newcommand{\eps}{\varepsilon}
\newcommand{\dd}{\,\mathrm{d}}
\begin{document}
\section*{UnsolvedMath ID: 701509}
\textbf{Source problem code:} catalogue entry
\subsection{Definitions and mathematical statement}
The supplied record does not explicitly specify the coefficient domain \(K\). Conditional on the usual interpretation that \(K\) is a field, the proposed statement is as follows.

Let \(K\) be any field and let \(G\neq\{1\}\) be any nontrivial group. The group algebra \(K[G]\) consists of finite formal sums \(\sum_{g\in G}a_g g\), with coefficients in \(K\), and multiplication induced by the group operation. Define the augmentation homomorphism and augmentation ideal by
\[
\varepsilon:K[G]\longrightarrow K,\qquad
\varepsilon\!\left(\sum_{g\in G}a_g g\right)=\sum_{g\in G}a_g,
\qquad I_K(G)=\ker\varepsilon.
\]
For a unital ring \(R\), let \(J(R)\) denote its Jacobson radical, the intersection of all maximal left ideals of \(R\). A group is locally finite if every finitely generated subgroup is finite. For a prime \(p\), a \(p\)-group here means a group in which every element has finite order equal to a power of \(p\).

The assertion attributed in the statement to I. Kaplansky is
\[
I_K(G)=J(K[G])
\quad\Longleftrightarrow\quad
\exists\text{ a prime }p\;\bigl(\operatorname{char}K=p>0
\ \text{and }G\text{ is a locally finite }p\text{-group}\bigr).
\]
The field hypothesis is an explicit interpretive qualification requiring confirmation from the original source; the supplied record does not determine a formulation for more general coefficient rings.
\subsection{Short English statement}
Under the explicitly conditional interpretation that \(K\) is a field, the question asks whether, for every nontrivial group \(G\), the augmentation ideal of \(K[G]\) equals its Jacobson radical exactly when \(K\) has positive prime characteristic \(p\) and every finitely generated subgroup of \(G\) is a finite \(p\)-group.
\subsection{Sources}
\begin{itemize}
\item Supplied local UnsolvedMath registry record, permanent ID \(701509\), exact source code \texttt{catalogue entry}, titled \emph{Kourovka Notebook Problem 18.80}; upstream revision \texttt{d371ee917ac9647fbb06cd6fc529deec6f3dda82}.
\item \emph{The Kourovka Notebook (Unsolved Problems in Group Theory)}, original problem 18.80, 18th issue (2014); the supplied background also identifies the 21st issue edition (2026), \texttt{arXiv:1401.0300v48}. The collection description names E. I. Khukhro and V. D. Mazurov as editors.
\item The record lists D. S. Passman as proposer, with proposed year 2014; the statement itself attributes the assertion to I. Kaplansky. Both attributions are retained.
\item The registry labels the problem \emph{open} and records an update date of 2026-10-16. These are supplied registry metadata, not an independently verified assessment of the literature. No separate literature-verification date or external source URL was supplied.
\end{itemize}
\end{document}
